\section{Problem Formulation}

Let $x_t$ denote the PM2.5 measurement at time $t$. For a forecast time $t$, the historical air-quality lookback window is
\begin{equation}
    X_t = [x_{t-L+1}, x_{t-L+2}, \ldots, x_t],
    \label{eq:lookback}
\end{equation}
where $L$ is the lookback length. Let $W_t$ denote weather covariates aligned
with the lookback window, $C_t$ denote calendar features, and $E_t$ denote event
context available under the recorded source assumption at the forecast origin.
The target associated with the input in (\ref{eq:lookback}) is
\begin{equation}
    Y_{t+1:t+H} = [x_{t+1}, x_{t+2}, \ldots, x_{t+H}],
    \label{eq:target}
\end{equation}
where $H$ is the forecasting horizon.
Equation~(\ref{eq:target}) therefore defines the 24-step vector evaluated at
every valid origin.

The forecasting goal is to learn a function
\begin{equation}
    \hat{Y}_{t+1:t+H} =
    F_{\theta}(X_t, W_t, C_t, E_t, R^{ts}_t, R^{ev}_t, D_t),
    \label{eq:forecast}
\end{equation}
where $R^{ts}_t$ is retrieved time-series knowledge, $R^{ev}_t$ is retrieved event knowledge, and $D_t$ contains concept-drift indicators. The implemented target setting is $L=168$ hours and $H=24$ hours.
Equation~(\ref{eq:forecast}) states the full information interface; the
implemented modules that construct these arguments are specified in the next
section.

The evaluation uses a time-based split:
\begin{equation}
    \mathcal{T}_{train} < \mathcal{T}_{val} < \mathcal{T}_{test},
    \label{eq:chronological_split}
\end{equation}
where the earliest 70\% of observations are used for training, the next 15\%
for validation, and the latest 15\% for testing. The ordering in
(\ref{eq:chronological_split}) is retained for feature fitting, retrieval
eligibility, model selection, and final evaluation.

Table~\ref{tab:symbols} summarizes the essential notation.

\begin{table}[!t]
\centering
\caption{Main Symbols Used in the Proposed Formulation}
\label{tab:symbols}
\begin{tabular}{ll}
\toprule
Symbol & Meaning \\
\midrule
$L$ & Lookback window length \\
$H$ & Forecast horizon \\
$X_t$ & Historical PM2.5 sequence at time $t$ \\
$W_t$ & Weather covariates \\
$C_t$ & Calendar/context features \\
$E_t$ & Event context available at the forecast origin \\
$R^{ts}_t$ & Retrieved time-series knowledge \\
$R^{ev}_t$ & Retrieved event knowledge \\
$D_t$ & Drift indicators \\
$q_t$ & SACB context vector in $\mathbb{R}^{86}$ \\
$G_t$ & LSER top-$k$ retrieved set \\
$\rho(G_t)$ & Retrieval summary in $\mathbb{R}^{51}$ \\
$\mathcal{O}_t$ & Machine-readable DFEH evidence record \\
$\hat{Y}_{t+1:t+H}$ & Predicted future PM2.5 values \\
\bottomrule
\end{tabular}
\end{table}
